\section{Literature review}

This chapter provides an overview of the existing research and establishes the theoretical foundation for the work presented in this thesis.

\subsection{Background}

The field under investigation has seen significant developments over the past decade.
Early approaches focused on manual processes, which were time-consuming and error-prone \cite{sample-design}.
As the field matured, automated solutions emerged that addressed many of these limitations.

Table~\ref{tab:comparison} presents a comparison of existing approaches.

\begin{table}[H]
    \caption{Comparison of existing approaches} \label{tab:comparison}
    \begin{tabular}{|l|c|c|c|}
    \hline \textbf{Approach} & \textbf{Accuracy (\%)} & \textbf{Speed (ms)} & \textbf{Scalability} \\
    \hline Method A & 85 & 120 & Low \\
    \hline Method B & 90 & 95  & Medium \\
    \hline Method C & 88 & 110 & High \\
    \hline Method D & 92 & 80  & Medium \\
    \hline
    \end{tabular}
\end{table}

\subsection{Theoretical framework}

The theoretical framework for this research draws from several disciplines, including computer science, statistics, and domain-specific knowledge.
The following equation describes the fundamental relationship:

\begin{equation} \label{eq:basic}
    E = \sum_{i=1}^{n} w_i \cdot f(x_i)
\end{equation}

where $E$ is the overall effectiveness metric, $w_i$ are the weight coefficients, and $f(x_i)$ represents the performance function for each component.

\subsection{Related work}

Several researchers have contributed to this field.
Smith and Johnson \cite{sample-ref} proposed an approach based on modular architecture that achieved promising results.
Brown et al. \cite{sample-analysis} extended this work by introducing optimization techniques that improved performance by approximately 15\%.
Lee and Park \cite{sample-testing} conducted a comprehensive evaluation of multiple methods and identified key factors affecting system performance.

\subsection{Standards and best practices}

Industry standards play an important role in ensuring quality and interoperability.
Organizations such as \hyperlink{ieee}{IEEE}, \hyperlink{iso}{ISO}, and \hyperlink{nist}{NIST} have published guidelines that inform the design of modern systems \cite{nist-guide}.

Figure~\ref{fig:architecture} illustrates a typical system architecture based on these standards.

\begin{figure}[H]
    \centering
    \begin{tikzpicture}[
        box/.style={rectangle, draw, minimum width=2.5cm, minimum height=0.8cm, align=center},
        >=Stealth
    ]
        \node[box] (ui) {User Interface};
        \node[box, below=0.6cm of ui] (api) {API Gateway};
        \node[box, below left=0.6cm and 0.3cm of api] (svc1) {Service A};
        \node[box, below right=0.6cm and 0.3cm of api] (svc2) {Service B};
        \node[box, below=0.6cm of api] (db) {Database};
        \draw[->] (ui) -- (api);
        \draw[->] (api) -- (svc1);
        \draw[->] (api) -- (svc2);
        \draw[->] (svc1) -- (db);
        \draw[->] (svc2) -- (db);
    \end{tikzpicture}
    \caption{General system architecture} \label{fig:architecture}
\end{figure}

\subsection{Gap analysis}

Despite significant progress, several gaps remain in the current body of knowledge:
\begin{itemize}
  \item existing methods do not adequately address scalability;
  \item performance optimization under resource constraints is underexplored;
  \item integration with modern deployment pipelines requires further investigation.
\end{itemize}

This thesis addresses these gaps by proposing a novel approach that combines the strengths of existing methods while mitigating their weaknesses.

\clearpage
